\chapter{The Scenario Authoring Language}
\label{ch:authoring-language}

\index{scenario}
\index{observable event}
This chapter proposes the first concrete surface language. It is intentionally a
small Lean DSL: ordinary \texttt{def}, function parameters, \texttt{let}, typed
values, and \texttt{do} notation provide the host language; a fixed library adds
scenario, endpoint, choice, probability, and rendering operations. The examples
are a design target. Names may change as the Lean implementation makes their
semantics precise.

\section{The smallest useful scenario}

A scenario is a named grammar fragment. Every emitted terminal is an observable
message event. Endpoint commands establish the \texttt{from} and \texttt{to}
context for the messages that follow.

\begin{lstlisting}
def requestQuote (customer shop : Actor) : Scenario :=
  scenario "request a quote" do
    let requestId <- declare RequestId "requestId"

    from_ customer
    to_ shop
    emit "quote requested" [field "requestId" requestId]

    let quoteId <- declare QuoteId "quoteId"
    let amount  <- declare Money "amount"

    from_ shop
    to_ customer
    emit "quote offered"
      [ field "requestId" requestId
      , field "quoteId" quoteId
      , field "amount" amount
      ]
\end{lstlisting}

The spelling \texttt{from\_} avoids collision with Lean syntax while remaining
visually explicit. Endpoint commands affect authoring context only; elaboration
copies both endpoints into every event. A scenario with an \texttt{emit} before
both endpoints have been established is rejected.

\subsection{Variables and assignment}

There are two kinds of assignment.

\begin{itemize}
  \item Ordinary \texttt{let x := expression} evaluates a known Lean expression.
  \item \texttt{let x <- declare T "x"} introduces a symbolic value of type
        \texttt{T}. The grammar binds it when generating or matching a trace, and
        later events may reuse it.
\end{itemize}

Variables are immutable. A later \texttt{let attempt := attempt + 1} creates a new
binding rather than mutable hidden state. Information flow is consequently visible
in the grammar's data dependencies. Equality of two occurrences of
\texttt{requestId} means they carry the same value, not merely values with the same
field name.

A value received from the environment is bound by matching an observable event:

\begin{lstlisting}
from_ customer
to_ shop
receive "quote accepted" do
  let quoteId <- take QuoteId "quoteId"
  let payment <- take PaymentToken "paymentToken"

from_ shop
to_ bank
emit "payment authorization requested"
  [field "quoteId" quoteId, field "paymentToken" payment]
\end{lstlisting}

\texttt{receive} and \texttt{emit} both denote grammar terminals; they differ only
in the viewpoint used when projecting a global scenario to executable participant
code.

\section{Vague messages first, typed alphabets later}

\index{alphabet!vague}
Early requirements should not require a finished protobuf schema. A vague message
has a stable semantic name and a list of named, typed fields. It is already precise
enough to express ordering, correlation, and temporal properties.

\begin{lstlisting}
emit "quote offered"
  [field "requestId" requestId, field "amount" amount]
\end{lstlisting}

Later, a refinement maps that terminal to a concrete alphabet:

\begin{lstlisting}
refineMessage "quote offered" with Proto.Shop.QuoteOffered where
  request_id := field "requestId"
  quote_id   := field "quoteId"
  cents      := encodeMoney (field "amount")
\end{lstlisting}

The refinement obligation is language inclusion: every concrete implementation
trace, after erasing representation details, must be admitted by the abstract
grammar. Refinement may decide serialization, retry mechanics, and internal state;
it must not invent externally visible behaviour forbidden by the requirements.

\section{One concern per scenario}

Scenarios should be small enough to read as stories and independent enough to
render separately. The quote protocol can be decomposed as follows.

\begin{lstlisting}
def quoteAccepted (customer shop : Actor) : Scenario :=
  scenario "quote accepted" do
    include (requestQuote customer shop)
    -- bind quoteId from the included scenario's exports
    use quoteId <- exported QuoteId "quoteId"
    from_ customer; to_ shop
    emit "quote accepted" [field "quoteId" quoteId]

def quoteDeclined (customer shop : Actor) : Scenario :=
  scenario "quote declined" do
    include (requestQuote customer shop)
    use quoteId <- exported QuoteId "quoteId"
    from_ customer; to_ shop
    emit "quote declined" [field "quoteId" quoteId]

def quoteOutcome (customer shop : Actor) : Scenario :=
  oneOf "customer decision"
    [ quoteAccepted customer shop
    , quoteDeclined customer shop
    ]
\end{lstlisting}

\texttt{include} is textual only in the rendered book view; semantically it is a
reference to a named nonterminal. The compiler reads all referenced scenarios,
checks exported bindings, detects cycles, and builds one grammar. Tools retain the
scenario boundary for navigation, diagrams, tests, and counterexamples.

\section{Interaction diagram: accepted quote}

The first rendering of any scenario is a sequence-style interaction diagram.
It shows observable messages and parameter flow without suggesting an internal
implementation architecture.

\begin{figure}[htbp]
\centering
\begin{tikzpicture}[>=Latex,actor/.style={draw,rounded corners,minimum width=2.6cm,
  minimum height=.7cm,fill=black!4},message/.style={->,thick}]
  \node[actor] (customer) at (0,0) {Customer};
  \node[actor] (shop) at (7,0) {Shop};
  \draw[dashed] (customer.south) -- ++(0,-5.2);
  \draw[dashed] (shop.south) -- ++(0,-5.2);
  \draw[message] (0,-1) -- node[above,sloped]{quote requested(requestId)} (7,-1);
  \draw[message] (7,-2) -- node[above,sloped]{quote offered(requestId, quoteId, amount)} (0,-2);
  \draw[message] (0,-3) -- node[above,sloped]{quote accepted(quoteId)} (7,-3);
  \draw[message] (7,-4) -- node[above,sloped]{order confirmed(quoteId, orderId)} (0,-4);
\end{tikzpicture}
\caption{Interaction diagram derived from the accepted-quote scenario.}
\label{fig:quote-interaction}
\end{figure}

\section{Choices and unresolved probability}

\index{probability}
Nondeterminism and probability are different. \texttt{oneOf} says that multiple
behaviours are legal without claiming how often they occur. \texttt{weighted}
attaches a distribution when the abstraction has one.

\begin{lstlisting}
def paymentOutcome : Scenario :=
  oneOf "payment outcome" [paymentApproved, paymentRejected]

def estimatedPaymentOutcome : Scenario :=
  choose "estimated payment outcome"
    [ weighted (95 / 100) paymentApproved
    , weighted ( 5 / 100) paymentRejected
    ]
\end{lstlisting}

Probabilities belong to choices, not to messages. An early requirement can use
\texttt{oneOf}, a symbolic weight, or a range such as \texttt{between 0.90 0.99}.
Implementation refinement resolves those claims to a well-defined stochastic
model or to deterministic conditions. Refinement must preserve possibility even
when a current implementation assigns probability zero to a legal abstract branch;
otherwise changing operational policy would silently change the contract.

\section{Composing requirements means intersecting constraints}

\index{composition!requirement intersection}
Scenario composition builds behaviour in familiar language-theoretic ways:

\begin{center}
\begin{tabular}{@{}lll@{}}
\toprule
Operator & Language operation & Purpose \\
\midrule
\texttt{then a b} & concatenation & $a$ happens before $b$ \\
\texttt{oneOf [a,b]} & union & either story is legal \\
\texttt{together a b} & shuffle/interleaving & independent concurrent stories \\
\texttt{repeat a} & closure & repeated conversation \\
\texttt{where\_ p a} & trace filtering & context-sensitive legality \\
\bottomrule
\end{tabular}
\end{center}

Combining requirements is subtler. A requirement usually does not add another
message sequence; it removes histories that violate a concern. Therefore
\texttt{requireAll} denotes intersection of accepted trace sets.

\begin{lstlisting}
def quoteSystem : SystemSpec :=
  generatedBy quoteOutcome
  |> requireAll
      [ correlate "offered quote matches request"
          (sameField "requestId")
      , forbid "confirmation without acceptance"
          (event "order confirmed" && notPreviously "quote accepted")
      , ctl "every accepted quote is answered"
          (AG (accepted --> AF (confirmed || rejected)))
      , timing "answer within service objective"
          (accepted --> within 2.seconds (confirmed || rejected))
      ]
\end{lstlisting}

This distinction prevents an important category error: union and concatenation
assemble possible stories; intersection layers independent obligations over those
stories. A system is consistent only if the intersection remains nonempty. When it
is empty, tooling should return a minimal conflicting set of requirements and the
shortest available counterexample prefix.

\section{State-machine rendering}

The same scenario grammar can be rendered as a state machine. States are grammar
residuals---the legal continuations after a trace prefix---rather than snapshots of
private implementation memory.

\begin{figure}[htbp]
\centering
\begin{tikzpicture}[->,>=Latex,node distance=1.8cm,on grid,
  every state/.style={draw,minimum size=1.05cm}]
  \node[state,initial] (start) {start};
  \node[state] (asked) [below=of start] {asked};
  \node[state] (offered) [below=of asked] {offered};
  \node[state,accepting] (accepted) [below left=1.8cm and 2.2cm of offered] {accepted};
  \node[state,accepting] (declined) [below right=1.8cm and 2.2cm of offered] {declined};
  \path
    (start) edge node[right,fill=white,inner sep=1pt] {quote requested} (asked)
    (asked) edge node[right,fill=white,inner sep=1pt] {quote offered} (offered)
    (offered) edge node[left,fill=white,inner sep=1pt,align=right] {quote accepted} (accepted)
    (offered) edge node[right,fill=white,inner sep=1pt,align=left] {quote declined} (declined);
\end{tikzpicture}
\caption{A state machine derived from grammar residuals.}
\label{fig:quote-state-machine}
\end{figure}

\section{Charts are declarative trace reductions}

\index{chart specification}
Charts are also specified in the Lean subset. A chart names an observable event
source, a grouping or $x$ expression, and an aggregate or $y$ expression. It does
not contain rendering code.

\begin{lstlisting}
def outcomeShare : ChartSpec :=
  pie "quote outcomes" do
    source events
    where_ (kind == "quote accepted" || kind == "quote declined")
    groupBy kind
    value count

def responseLatency : ChartSpec :=
  line "quote response latency" do
    source (correlate events by "quoteId")
    x tstamp
    y (elapsed "quote offered" "customer decision")
    series "p50" median
\end{lstlisting}

\begin{figure}[htbp]
\centering
\begin{minipage}{.43\textwidth}
\centering
\begin{tikzpicture}
  \def\r{1.55}
  \fill[leanblue!80] (0,0) -- (90:\r) arc (90:-234:\r) -- cycle;
  \fill[orange!80] (0,0) -- (-234:\r) arc (-234:-270:\r) -- cycle;
  \node at (0,-2.05) {accepted 90\%};
  \node at (0,-2.45) {declined 10\%};
\end{tikzpicture}
\end{minipage}\hfill
\begin{minipage}{.52\textwidth}
\centering
\resizebox{.97\linewidth}{!}{%
\begin{tikzpicture}
\begin{axis}[
  width=\linewidth,height=5cm,
  xlabel={event time},ylabel={latency (ms)},
  xmin=1,xmax=5,ymin=0,ymax=500,
  grid=major]
  \addplot[color=leanblue,mark=*] coordinates {(1,310) (2,240) (3,180) (4,205) (5,145)};
\end{axis}
\end{tikzpicture}
}
\end{minipage}
\caption{Example pie and line renderings from declarative chart specifications.}
\label{fig:charts}
\end{figure}

Interaction diagrams, state machines, and charts are projections of the same typed
scenario and event definitions. Keeping them derived prevents a diagram from
becoming a second, drifting specification.
